% concatenated from main.tex
\documentclass[table]{amsart}
\usepackage[margin=1.1in]{geometry}
\usepackage{amscd,amsmath,amsxtra,amsthm,amssymb,stmaryrd,xr,mathrsfs,mathtools,enumerate,commath, comment, mathtools, orcidlink}
\usepackage{tikz}
\usetikzlibrary{arrows.meta,positioning}

\usetikzlibrary{calc}
\usepackage{tikz-3dplot}
\usepackage{stmaryrd}
\usepackage{multirow}
\usepackage{xcolor}
\usepackage{commath}
\usepackage{comment}
\usepackage{svg}
\usepackage{graphics}
%\usepackage[anticlockwise]{rotating}
\usepackage{longtable} % for 'longtable' environment
\usepackage{pdflscape} % for 'landscape' environment
\usepackage{booktabs}
\usepackage{hyperref}
\definecolor{vegasgold}{rgb}{0.77, 0.7, 0.35}
\definecolor{darkgoldenrod}{rgb}{0.72, 0.53, 0.04}
\definecolor{gold(metallic)}{rgb}{0.83, 0.69, 0.22}
\hypersetup{
 colorlinks=true,
 linkcolor=darkgoldenrod,
 filecolor=brown,      
 urlcolor=gold(metallic),
 citecolor=darkgoldenrod,
 }
\newtheorem{lthm}{Theorem}
\renewcommand{\thelthm}{\Alph{lthm}}
\newtheorem{lcor}{Corollary}
\renewcommand{\thelcor}{\Alph{lcor}}
\usepackage[all,cmtip]{xy}


\DeclareFontFamily{U}{wncy}{}
\DeclareFontShape{U}{wncy}{m}{n}{<->wncyr10}{}
\DeclareSymbolFont{mcy}{U}{wncy}{m}{n}
\DeclareMathSymbol{\Sh}{\mathord}{mcy}{"58}
\usepackage[T2A,T1]{fontenc}
\usepackage[OT2,T1]{fontenc}




\newtheorem{theorem}{Theorem}[section]
\newtheorem{lemma}[theorem]{Lemma}
\newtheorem{ass}[theorem]{Assumption}
\newtheorem*{theorem*}{Theorem}
\newtheorem*{ass*}{Assumption}
\newtheorem{definition}[theorem]{Definition}
\newtheorem{corollary}[theorem]{Corollary}
\newtheorem{remark}[theorem]{Remark}
\newtheorem{example}[theorem]{Example}
\newtheorem{conjecture}[theorem]{Conjecture}
\newtheorem{proposition}[theorem]{Proposition}
\newtheorem{question}[theorem]{Question}

\newcommand{\tr}{\operatorname{tr}}
\newcommand{\Fil}{\operatorname{Fil}}
\newcommand{\cor}{\operatorname{cor}}
\newcommand{\ac}{\operatorname{ac}}
\newcommand{\ns}{\operatorname{ns}}
\newcommand{\bK}{\mathbb{K}}
\newcommand{\DD}{\mathbb{D}}
\newcommand{\BB}{\mathbb{B}}
\newcommand{\NN}{\mathbb{N}}
\newcommand{\bA}{\mathbf{A}}
\newcommand{\fE}{\mathfrak{E}}
\newcommand{\cT}{\mathcal{T}_1}
\newcommand{\bE}{\textbf{E}}
\newcommand{\cK}{\mathcal{K}}
\newcommand{\cP}{\mathcal{P}}
\newcommand{\cE}{\mathcal{E}}
\newcommand{\fL}{\mathfrak{L}}
\newcommand{\cI}{\mathcal{I}}
\newcommand{\fA}{\mathfrak{A}}
\newcommand{\Qp}{\mathbb{Q}_p}
\newcommand{\Zp}{\mathbb{Z}_p}
\newcommand{\pdiv}{\mid\!\mid}
\newcommand{\HH}{\mathbb{H}}
\newcommand{\FF}{\mathcal{F}}
\newcommand{\cG}{\mathcal{G}}
\newcommand{\Kant}{K_{\operatorname{anti}}}
\newcommand{\sH}{\mathscr{H}}
\newcommand{\uk}{\underline{k}}
\newcommand{\Frac}{\mathrm{Frac}}
\newcommand{\ord}{\mathrm{ord}}
\newcommand{\bG}{\mathbf{G}}
\newcommand{\bT}{\mathbf{T}}
\newcommand{\bB}{\mathbf{B}}
\newcommand{\bU}{\mathbf{U}}
\newcommand{\cH}{\mathcal{H}}
\newcommand{\cris}{\mathrm{cris}}
\newcommand{\I}{\op{I}_n}
\newcommand{\lb}{[\![}
\newcommand{\rb}{]\!]}
\newcommand{\rhobar}{\barrho}

\newcommand{\Z}{\mathbb{Z}}
\newcommand{\Q}{\mathbb{Q}}
\newcommand{\F}{\mathbb{F}}
\newcommand{\D}{\mathfrak{D}}
\newcommand{\fM}{\mathfrak{M}}
\newcommand{\Lp}{\mathbf{L}_p}
\newcommand{\vp}{\varphi}
\newcommand{\cC}{\mathcal{C}}
\newcommand{\cO}{\mathcal{O}}
\newcommand{\cS}{\mathcal{S}}
\newcommand{\Iw}{\mathrm{Iw}}
\newcommand{\HIw}{H^1_{\mathrm{Iw}}}
\newcommand{\rk}{\mathrm{rk}}
\newcommand{\fv}{\mathfrak{v}}
\newcommand{\n}{\mathfrak{n}}

\newcommand{\loc}{\mathrm{loc}}
\newcommand{\cc}{\mathfrak{c}}
\newcommand{\ff}{\mathfrak{f}}
\newcommand{\fI}{\mathfrak{I}}
\newcommand{\Qpn}{\QQ_{p,n}}
\newcommand{\Hom}{\mathrm{Hom}}
\newcommand{\new}{\mathrm{new}}
\newcommand{\base}{\mathrm{base}}
\newcommand{\Sel}{\mathrm{Sel}}
%\newcommand{\mathbb{C}har}{\mathrm{char}}
\newcommand{\Spec}{\mathrm{Spec}\ }
\newcommand{\Gr}{\mathrm{Gr}}
\newcommand{\Reg}{\mathrm{Reg}}
\newcommand{\len}{\mathrm{len}}
\newcommand{\coker}{\mathrm{coker}}
\newcommand{\pr}{\mathrm{pr}}
\newcommand{\Tw}{\mathrm{Tw}}
\newcommand{\rank}{\mathrm{rank}}
\renewcommand{\div}{\mathrm{div}}
\newcommand{\corank}{\mathrm{corank}}
\newcommand{\ur}{\mathrm{ur}}
\newcommand{\N}{\mathbb{N}}
\newcommand{\Qinf}{\QQ_\infty}
\newcommand{\Qpinf}{\QQ_{\infty,p}}
\newcommand{\Qvinf}{\QQ_{\infty,v}}
\newcommand{\Kinf}{K_\infty}
\newcommand{\Kvinf}{K_{\infty,v}}
\newcommand{\Kpinf}{K_{\infty,p}}
\newcommand{\vpi}{p}
\newcommand{\Afs}{A_{f,s}}
\newcommand{\Ags}{A_{g,s}}
\newcommand{\Tfs}{\T_{f,s}}
\newcommand{\Vfs}{\mathcal{V}_{f,s}}
\newcommand{\As}{\mathcal{A}_{f,s}}
\newcommand{\Kwinf}{K_{\infty,w}}
\newcommand{\Klinf}{_{\infty,\ell}}
\newcommand{\TT}{\mathbf{T}}
\renewcommand{\AA}{\mathbf{A}}
\newcommand{\X}{\mathfrak{X}}
\newcommand{\tor}{\mathrm{tors}}
\newcommand{\cyc}{\mathrm{cyc}}
\newcommand{\fm}{\mathfrak{m}}
\newcommand{\A}{\mathcal{A}}
\newcommand{\V}{\mathcal{V}}
\newcommand{\Hc}{\mathcal{H}}
\newcommand{\W}{\mathcal{W}}
\newcommand{\m}{\mathrm{max}}
\newcommand{\Prob}{\mathbb{P}}
\newcommand{\lk}{\operatorname{lk}}
\newcommand{\id}{\mathrm{id}}
\newcommand{\cover}{\Sigma_2}
\newcommand{\lkmod}{\overline{\lk}}

\renewcommand{\labelitemii}{$\diamond$}
\newcommand{\mb}[1]{\mathbf{#1}}
\newcommand{\green}[1]{\textcolor{green}{#1}}
\newcommand{\blue}[1]{\textcolor{blue}{#1}}
\newcommand{\red}[1]{\textcolor{red}{#1}}

\newcommand{\AR}[1]{\textcolor{blue}{#1}}
\newcommand{\TS}[1]{\textcolor{red}{#1}}

\newcommand{\op}[1]{\operatorname{#1}}
%\newcommand{\mathbb{C}}{\mathbb{C}}
\newcommand\mtx[4] { \left( {\begin{array}{cc}
 #1 & #2 \\
 #3 & #4 \\
 \end{array} } \right)}

 
 \newcommand{\widebar}[1]{\mkern 2.5mu\overline{\mkern-2.5mu#1\mkern-2.5mu}\mkern 2.5mu}
 
 \newcommand{\barrho}{\overline{\rho}}
 \DeclareMathSymbol{\sha}{\mathord}{mcy}{"58}
 \makeatletter
\newcommand{\mylabel}[2]{#2\def\@currentlabel{#2}\label{#1}}
\makeatother

\numberwithin{equation}{section}


\title[class groups of cubic orders with prescribed shape]{On the average $2$-torsion in class groups and narrow class groups of cubic orders with prescribed shape}
\author[A.~Ray]{Anwesh Ray\, \orcidlink{0000-0001-6946-1559}}
\address[Ray]{Chennai Mathematical Institute, H1, SIPCOT IT Park, Kelambakkam, Siruseri, Tamil Nadu 603103, India}
\email{anwesh@cmi.ac.in}

\begin{document}

\maketitle


\begin{abstract}
We study the distribution of $2$-torsion in class groups and narrow class groups
of cubic fields and cubic orders subject to prescribed shape conditions. The
\emph{shape} of a cubic order in a number field is a natural geometric invariant that takes values in the modular
surface $\mathbb{H}/\operatorname{GL}_2(\mathbb{Z})$. Fix a subset $W$ of the modular surface with positive hyperbolic measure whose boundary has measure $0$. Refining the methods of Bhargava and Varma, we prove that
among cubic fields with shape in $W$, the average size of the $2$-torsion subgroup
of the class group is $5/4$ for totally real fields and $3/2$ for complex fields,
while the average size of the $2$-torsion subgroup of the narrow class group for
totally real cubic fields is $2$. Similar results are obtained for cubic orders satisfying prescribed local conditions at all primes.
\end{abstract}


\section{Introduction}
\subsection{Motivation and historical context} The arithmetic statistics of number fields have advanced significantly in recent
years, largely motivated by the Cohen--Lenstra heuristics, which predict the
distribution of $p$-parts of class groups of number fields of fixed degree,
ordered by discriminant. Till date the only cases for which these conjectures are proven is when $p=3$ and $K$ varies over quadratic number fields and when $p=2$ and $K$ varies over cubic number fields. Davenport and Heilbronn \cite{davenportheilbronn} show that as $K$ varies over real (resp. imaginary) quadratic fields, the average size of $\op{Cl}_3(K)$ is $4/3$ (resp. $2$). On the other hand, Bhargava \cite{bhargavequartic} showed for totally real (resp. complex) cubic number fields $K$, the average size of $\op{Cl}_2(K)$ is $5/4$ (resp. $3/2$). These results are driven in large part by the development of powerful parametrization techniques originating in the work of Delone–Faddeev \cite{DeloneFaddev} and their far-reaching generalizations due to Bhargava \cite{bhargavahighercomplaws2}.
\par Alongside these developments, there has been increasing interest in finer
geometric invariants of number fields. One such invariant is the \emph{shape} of
a number field, defined using the Minkowski embedding of its ring of integers.
Equipping $K\otimes_{\mathbb{Q}}\mathbb{R}$ with the trace form, the image of
$\mathcal{O}_K$ becomes a lattice containing the vector corresponding to $1$.
Projecting orthogonally to this vector yields a rank-$(n-1)$ lattice, whose
equivalence class up to scaling and orthogonal transformations defines the
shape. The space of all such shapes may be identified with the double coset
space
\[
\mathcal{S}_{n-1}
=
\operatorname{GL}_{n-1}(\mathbb{Z})
\backslash
\operatorname{GL}_{n-1}(\mathbb{R})
/
\operatorname{GO}_{n-1}(\mathbb{R}),
\]
and carries a natural measure induced from Haar measure. Distribution questions for shapes were initiated by Terr \cite{Terr97}, who proved that the
shapes of cubic fields are equidistributed in $\mathcal{S}_2$. In this case,
$\mathcal{S}_2$ may be identified with the modular surface
$\mathbb{H}/\op{GL}_2(\mathbb{Z})$, where $\mathbb{H}$ is the complex upper half plane. Equidistribution holds
with respect to the hyperbolic measure $y^{-2}dxdy$. More recently, Bhargava and Harron \cite{bhargavaharron}
established analogous equidistribution results for $S_n$-number fields of
degrees $n=4,5$.
\subsection{Main results}
In this work, we combine these perspectives by studying averages of $2$-torsion
in class groups and narrow class groups of cubic orders subject to prescribed
shape conditions. We show that imposing such geometric restrictions does not
alter the limiting averages. The expected size of the relevant $2$-torsion
groups of cubic fields with shape contained in a suitably nice subset $W\subseteq \mathcal{S}_2$ remains coincides with the averages obtained by Bhargava and Varma \cite{BhargavaVarma}. The following is a refinement of \cite[Theorem 1]{BhargavaVarma}.
\begin{lthm}\label{thm a}
    Let $W \subset \mathcal{S}_2$ be a measurable subset with positive measure whose boundary is of measure
zero. Fix any finite set of primes and impose arbitrary
local conditions at these primes. When cubic fields with shape lying in $W$ and
satisfying these local conditions are ordered by absolute discriminant, the
following statements hold:
\begin{itemize}
\item[(a)] Among totally real cubic fields, the average size of the
$2$-torsion subgroup of the class group is $5/4$.
\item[(b)] Among complex cubic fields, the average size of the
$2$-torsion subgroup of the class group is $3/2$.
\item[(c)] Among totally real cubic fields, the average size of the
$2$-torsion subgroup of the narrow class group is $2$.
\end{itemize}

\end{lthm}

We also obtain results for all cubic orders. Let $(\Sigma_p)$ be an \emph{acceptable collection} of local specifications as in Definition \ref{acceptable Sigma}. Given an order $\cO$ in a cubic number field $K$, $\mathcal{I}(\cO)$ be the group of invertible fractional ideals of $\cO$ and $\mathcal{P}(\mathcal{O})\subseteq \mathcal{I}(\cO)$ be the subgroup of principal fractional ideals. Denote by $\mathcal{P}^+(\cO)
\subseteq \mathcal{P}(\mathcal{O})$ the subgroup of totally positive principal fractional ideals. Then the class group (resp. narrow class group) of $\cO$ is the quotient $\op{Cl}(\cO):=\mathcal{I}(\cO)/\mathcal{P}(\mathcal{O})$ (resp. $\op{Cl}^+(\cO):=\mathcal{I}(\cO)/\mathcal{P}^+(\mathcal{O})$). Given a natural number $n$, denote by $\op{Cl}_n(\cO)$ (resp. $\mathcal{I}_n(\cO)$) the $n$-torsion subgroup of $\op{Cl}(\cO)$ (resp. $\mathcal{I}(\cO)$). We have the following refinement of \cite[Theorem 2]{BhargavaVarma}.

\begin{lthm}\label{thm b}
    Let $(\Sigma_p)$ be an acceptable collection of local specifications, and let
$\Sigma$ denote the set of isomorphism classes of cubic orders
$\mathcal{O}$ such that $\mathcal{O}\otimes\mathbb{Z}_p \in \Sigma_p$ for all
primes $p$. Let $W \subset \mathcal{S}_2$ be a measurable subset with positive measure and whose boundary has
measure zero. When cubic orders $\mathcal{O}\in\Sigma$ whose shapes
lie in $W$ are ordered by absolute discriminant, the following hold:
\begin{itemize}
\item[(a)] For totally real cubic orders $\mathcal{O}$, the average value of
\[
\left|\mathrm{Cl}_{2}(\mathcal{O})\right|
-\tfrac{1}{4}\left|\mathcal{I}_{2}(\mathcal{O})\right|
\]
is $1$.
\item[(b)] For complex cubic orders $\mathcal{O}$, the average value of
\[
\left|\mathrm{Cl}_{2}(\mathcal{O})\right|
-\tfrac{1}{2}\left|\mathcal{I}_{2}(\mathcal{O})\right|
=
\left|\mathrm{Cl}_{2}^{+}(\mathcal{O})\right|
-\tfrac{1}{2}\left|\mathcal{I}_{2}(\mathcal{O})\right|
\]
is $1$.
\item[(c)] For totally real cubic orders $\mathcal{O}$, the average value of
\[
\left|\mathrm{Cl}_{2}^{+}(\mathcal{O})\right|
-\left|\mathcal{I}_{2}(\mathcal{O})\right|
\]
is $1$.
\end{itemize}

\end{lthm}
\noindent We note that Theorem \ref{thm b} is simply a restatement of Theorem \ref{main thm of article}. These results have many interesting consequences for cubic fields with precibed shape in a set $W\subseteq \mathcal{S}_2$. For instance, the following is a refinement of \cite[Corollary 3]{bhargavaharron}.
\begin{lthm}\label{thm c}
    Let $W \subset \mathcal{S}_2$ be a measurable subset with positive measure and whose boundary has
measure zero. When cubic fields with shape lying in $W$ and satisfying these local conditions
are ordered by absolute discriminant, the following statements hold:
\begin{itemize}
\item[(1)] A positive proportion of totally real cubic fields have odd class
number; in fact, this proportion is at least $75\%$.
\item[(2)] A positive proportion of complex cubic fields have odd class number;
in fact, this proportion is at least $50\%$.
\end{itemize}
\end{lthm}

\subsection{Organization}
The structure of the paper is as follows. In Section~2 we recall the notion of the shape of a lattice and of a number field, with particular emphasis on the cubic case. We also review the parametrizations of cubic rings and $2$-torsion ideal classes that will be used throughout the paper. In Section~3 we carry out the necessary volume computations, adapting the framework in \cite{bhargavaharron}. These calculations differ from those in \emph{loc. cit.} where the shape distribution is studied for quartic orders, while in our applications, the parameterization for $2$-torsion ideal classes in cubic orders is considered. The emphasis lies in the $\op{GL}_2$-action on the homogenous space $\mathbb{R}^2\otimes\op{Sym}^2(\mathbb{R}^3)$ as opposed to the $\op{GL}_3$-action. In Section~4 we incorporate local congruence conditions and establish equidistribution of shapes for acceptable families of cubic orders. Finally, we combine these results to prove our main theorems on the average size of $2$-torsion in class groups and narrow class groups under shape restrictions. Theorems \ref{thm a} and \ref{thm c} are proven at the end of this section.

\par We hope that these results further illustrate the robustness of arithmetic statistics under the imposition of natural geometric conditions, and provide additional evidence for the compatibility between algebraic and geometric aspects of number fields. 


\section{Preliminary notions}

\subsection{Shapes of lattices}

The notion of the \emph{shape} of a full rank lattice in Euclidean space is an invariant which captures its intrinsic geometry up to the natural symmetries. To formalize this notion, fix a real vector space $V$ of dimension $r$ with ordered basis $e_1, \dots, e_r$. Assume that $V$ is equipped with the standard inner product defined by $\langle e_i, e_j\rangle =\delta_{i,j}$. The \emph{space of shapes of rank-$r$ lattices} is defined to be the double coset space:
\[
\mathcal{S}_r := \op{GL}_r(\mathbb{Z}) \backslash \op{GL}_r(\mathbb{R}) / \op{GO}_r(\mathbb{R}),
\]
where $\op{GO}_r(\mathbb{R})$ denotes the group of invertible linear transformations preserving $\langle \cdot, \cdot \rangle$ up to a positive scalar multiple. Let $\Lambda\subset V$ be a full rank lattice in $V$ and let $\{b_1,\dots,b_r\}$ be a chosen $\Z$-basis of $\Lambda$. We write $b_i=\sum_{j=1}^n m_{i,j} e_j$ and let $M:=(m_{i,j})\in\op{GL}_r(\mathbb{R})$. Then the class of $M$ in the quotient $\op{GL}_r(\mathbb{Z}) \backslash \op{GL}_r(\mathbb{R})$ depends only on $\Lambda$ and not on the chosen basis. The \emph{shape of $\Lambda$} is defined to be the class of $M$ in the double coset space $\mathcal{S}_r$ and is denoted by $\op{sh}(\Lambda)$. The space of shapes carries a natural measure $\mu_{\mathcal{S}_r}$ induced from Haar measure on $\op{GL}_r(\mathbb{R})$.

\par An equivalent description is obtained via Gram matrices. Let $\mathcal{P}_r$ denote the space of positive definite symmetric $r\times r$ real matrices. The group $\op{GL}_r(\mathbb{Z})$ acts on $\mathcal{P}_r$ by $M\cdot G :=M^{T} G M$ for $G\in \mathcal{P}_r$ and $M\in \op{GL}_r(\mathbb{Z})$. On the other hand,
while $\lambda\in \mathbb{R}^{\times}$ acts by $\lambda\cdot G:=\lambda^2G$. Writing
\[
\mathcal{H}_r := \mathcal{P}_r / \mathbb{R}^{\times}
 = \{\, G \in M_r(\mathbb{R}) \mid G = G^T,\; G>0,\; \det G = 1 \,\},
\]
there is a natural identification
\[
\op{GL}_r(\mathbb{R}) / \op{GO}_r(\mathbb{R}) \;\xrightarrow{\;\sim\;}\; \mathcal{H}_r,\]
mapping $M$ to $|\det M|^{-2/r} M^T M$. Thus we have an isomorphism
\[
\mathcal{S}_r \simeq \op{GL}_r(\mathbb{Z}) \backslash \mathcal{H}_r.
\]
Under this identification, the shape of a lattice $\Lambda$ is represented by the class of its Gram matrix defined by $\op{Gr}(\Lambda):=\left(\langle v_i, v_j\rangle \right)$, and the relation
\[
\op{Gr}(\lambda A\cdot \Lambda)=\lambda^2 A^T \op{Gr}(\Lambda) A
\]
shows that this class is independent of the choice of basis and of scaling.

\par The case $r=2$ admits a particularly transparent description. Writing a general element of $\mathcal{H}_2$ as $\begin{pmatrix} a & b \\ b & c \end{pmatrix}$ where $a>0$ and $ac-b^2=1$, one obtains an explicit $\op{GL}_2(\mathbb{R})$-equivariant identification of $\mathcal{H}_2$ with $\mathbb{H}$, the complex upper half-plane. Explicitly, the matrix above corresponds to the point $z=x+iy\in\mathbb{H}$ with
\[
x=\frac{b}{c}, \qquad y=\sqrt{\frac{a}{c}-x^2}.
\]
Under this correspondence, the action of $\op{GL}_2(\mathbb{Z})$ on $\mathcal{H}_2$ becomes the usual fractional linear action on $\mathbb{H}$, and a fundamental domain is given by
\[
\mathcal{F}_2=\{\, x+iy\in\mathbb{H} \mid x\in[0,1/2),\; y>0,\; x^2+y^2\ge 1 \,\}.
\]
The induced measure on $\mathcal{F}_2$ is the hyperbolic measure $\frac{dx\, dy}{y^2}$ coming from Haar measure on $\op{GL}_2(\mathbb{R})$.
\par Let $K$ be a number field of degree $n$ with ring of integers $\mathcal{O}_K$. Let $ \{1, \alpha_1, \dots, \alpha_{n-1}\} $ be an integral basis and let $ \{\sigma_i\}_{1 \leq i \leq n} $ denote the set of embeddings of $ K $ into $ \mathbb{C}$.  The Minkowski embedding  
\[
\begin{aligned}
    j: K &\to \mathbb{C}^{n} \\
    \alpha &\mapsto \big( \sigma_1(\alpha), \dots, \sigma_n(\alpha) \big)
\end{aligned}
\]
identifies $ K $ with an $ n $-dimensional real subspace of $ \mathbb{C}^n $, which we denote by $K_{\mathbb{R}}$. This real vector space is naturally endowed with the inner product induced by the trace form $\langle x,y\rangle=\op{Tr}_{K/\mathbb{Q}}(xy)$, and hence $j(\mathcal{O}_K)$ acquires a well-defined Gram matrix. Since the vector $j(1)$ spans a distinguished one-dimensional subspace, it is customary to define the \emph{shape of $K$} to be the shape of the orthogonal projection of $j(\mathcal{O}_K)$ onto the hyperplane perpendicular to $j(1)$. This yields a point of $\mathcal{S}_{n-1}$. More concretely, let $\cO_K^\perp$ denote the image of the map $\alpha^{\perp} := n\alpha - \operatorname{tr}(\alpha)$. The \emph{shape of $K$}, denoted $\op{sh}_K$, is defined to be the shape of the lattice $j(\cO_K^\perp)$. It is represented by the Gram-matrix, which is given by
Gram matrix is given as follows: 
\[\label{gram matrix computation}
\operatorname{Gr}\left(j\left(\mathcal{O}_{K}^{\perp}\right)\right) = \big( \langle j(\alpha_i^{\perp}), j(\alpha_j^{\perp}) \rangle \big).\]

For quadratic fields, this invariant is trivial, as all rank-$1$ lattices are homothetic. In contrast, for cubic fields the associated projected lattice has rank $2$, and hence the shape of a cubic field is represented by a point in the fundamental domain $\mathcal{F}_2 \subset \mathbb{H}$.
\subsection{Parameterization of order $2$ ideals in cubic orders}

Let $T$ be a principal ideal domain. A ring of rank $n$ over $T$ is defined to be commutative $T$-algebra for which the underlying $T$-module $(R,+)$ is isomorphic to the free module $T^n$. When $n=2,3,4$ or $5$ we say that $R$ is a quadratic, cubic, quartic or quintic ring over $T$ respectively. A ring $R$ of rank $n$ is said to be \emph{nondegenerate} if $\op{Disc}(R)\neq 0$. 

\par We first recall the classical parametrization of cubic rings due to Delone and Faddeev \cite{DeloneFaddev}, as later extended and reformulated by Gan, Gross, and Savin \cite[Proposition 4.2]{GGS}. Let $f(x,y)=ax^{3}+bx^{2}y+cxy^{2}+dy^{3}$ be a binary cubic form with coefficients in $T$ and let $\gamma\in \op{GL}_2(T)$. Then \(\op{GL}_{2}(T)\) acts on the space of binary cubic forms via the twisted action:
\[
(\gamma\cdot f)(x,y)=\det(\gamma)^{-1} f\big((x,y)\gamma\big).
\]

\begin{theorem}[Delone--Faddeev, Gan--Gross--Savin]
There is a canonical bijection between isomorphism classes of cubic rings $R$ over $T$ and
$\op{GL}_{2}(T)$-equivalence classes of binary cubic forms $f$ with coefficients in $T$.
Under this correspondence, let $R(f)$ be the cubic ring associated to a form $f$. The discriminant of $R(f)$ coincides with the
discriminant of $f$. In particular, \(R(f)\) is nondegenerate if and only if \(\op{Disc}(f)\neq 0\).
\end{theorem}

\begin{proof} We provide a sketch of the proof for the benefit of exposition. To describe the correspondence explicitly, let \(R\) be a cubic ring. Choosing a \(T\)-basis \(\langle 1,\omega,\theta\rangle\) for \(R\), and translating \(\omega\) and \(\theta\) by suitable integers, we may assume that \(\omega\theta\in T\). Such a basis is called \emph{normal}. With respect to a normal basis, multiplication in \(R\) is given by relations of the form
\[
\omega\theta = n, \qquad
\omega^{2} = m + b\omega - a\theta, \qquad
\theta^{2} = \ell + d\omega - c\theta,
\]
for integers \(a,b,c,d,\ell,m,n\). Associativity of multiplication imposes constraints on \(\ell,m,n\), and a direct calculation shows that these constants are uniquely determined by
\[
n=-ad, \qquad m=-ac, \qquad \ell=-bd.
\]
The coefficients \(a,b,c,d\) therefore determine the cubic ring \(R\) uniquely, and one associates to \(R\) the binary cubic form
\[
f(x,y)=ax^{3}+bx^{2}y+cxy^{2}+dy^{3}.
\]

Conversely, starting from a binary cubic form \(f(x,y)=ax^{3}+bx^{2}y+cxy^{2}+dy^{3}\), the above relations define a cubic ring \(R(f)\) with underlying \(T\)-module \(T\oplus T\omega\oplus T\theta\). The associativity relations ensure that this construction is well-defined and inverse to the assignment \(R\mapsto f\). A direct computation shows that the discriminant of the cubic ring \(R(f)\) coincides with the discriminant of the associated binary cubic form,
\[
\op{Disc}(R(f))=\op{Disc}(f)
= b^{2}c^{2}-4ac^{3}-4b^{3}d-27a^{2}d^{2}+18abcd.
\]
\end{proof}
\par
A key input in our arguments is Bhargava’s parametrization
\cite[Theorem~4]{bhargavahighercomplaws2} of ideal classes of order two in cubic
rings via orbits of pairs of ternary quadratic forms. We briefly recall the aspects of this
parametrization that are relevant for our work, following the notation of
\cite{BhargavaVarma}. Let $V_T := T^{2}\otimes \op{Sym}^{2} T^{3}$ be
the space of pairs of ternary quadratic forms over \(T\). Equivalently, \(V_T\)
may be identified with the space of pairs of symmetric \(3\times 3\) matrices
with entries in \(T\),
\[
(A,B)=
\left(
\begin{bmatrix}
a_{11} & a_{12} & a_{13}\\
a_{12} & a_{22} & a_{23}\\
a_{13} & a_{23} & a_{33}
\end{bmatrix},
\begin{bmatrix}
b_{11} & b_{12} & b_{13}\\
b_{12} & b_{22} & b_{23}\\
b_{13} & b_{23} & b_{33}
\end{bmatrix}
\right).
\]
\noindent There is a natural algebraic action of the group
\(G_T:=\op{GL}_2(T)\times \op{GL}_3(T)\) on \(V_T\). Explicitly, for
\(\gamma=\mtx{r}{s}{t}{u}\in \op{GL}_2(T)\), define
\[
\gamma^{*}:=\mtx{r}{-t}{-s}{u}.
\]
Then the action of \(\op{GL}_2(T)\) on \(V_T\) is given by
\[
\gamma\cdot (A,B) := (A,B)\gamma^{*} = (rA - sB,\,-tA + uB).
\]
On the other hand, the action of \(\op{GL}_3(T)\) is induced from its natural
action on the underlying module \(T^{3}\): for \(g\in \op{GL}_3(T)\),
\[
g\cdot (A,B) := (gAg^{t},\, gBg^{t}).
\]

These two actions are compatible, yielding a joint action of
\(\op{GL}_2(T)\times \op{GL}_3(T)\) on \(V_T\). Let \(\Delta\) denote the subgroup
of \(\op{GL}_2(T)\times \op{GL}_3(T)\) consisting of elements of the form
\[
T_\lambda :=
\left(
\begin{bmatrix}
\lambda^{-2} & 0\\
0 & \lambda^{-2}
\end{bmatrix},
\begin{bmatrix}
\lambda & 0 & 0\\
0 & \lambda & 0\\
0 & 0 & \lambda
\end{bmatrix}
\right),
\qquad \lambda\in T^{\times}.
\]
A direct verification shows that \(T_\lambda\) acts trivially on \(V_T\). It
follows that the action of \(\op{GL}_2(T)\times \op{GL}_3(T)\) descends to an
action of the quotient group $G_T/\Delta$. To any pair \((A,B)\in V_T\), one associates a binary cubic form $f_{(A,B)}(x,y) := \op{Det}(Ax - By)$.
The coefficients of \(f_{(A,B)}\) are polynomial invariants under the
\(\op{SL}_3\)-action, and the discriminant
\[
\op{Disc}(A,B) := \op{Disc}\bigl(f_{(A,B)}\bigr)
\]
is a homogeneous polynomial of degree \(12\) in the coefficients of \(A\) and
\(B\). This discriminant generates the ring of relative invariants for the
action of \(\op{SL}_2\times \op{SL}_3\). An element \((A,B)\in V_T\) is called
\emph{nondegenerate} if \(\op{Disc}(A,B)\neq 0\).


\par
Via the Delone--Faddeev--Gan--Gross--Savin correspondence, the binary cubic form
associated to a pair of ternary quadratic forms determines a cubic ring over
\(T\). This relationship admits a substantial refinement that incorporates ideal
classes of order two. To state it precisely, fix a set \(S\) of representatives
for the left action of \(T^{\times}\) on \(T\setminus\{0\}\). For an ideal \(I\) of a
cubic ring over \(T\), we denote by \(N_{S}(I)\) the distinguished generator in
\(S\) of the ideal norm of \(I\) viewed as a \(T\)-ideal. 

\par Let $\mathcal{S}_T$ be the set of equivalence classes of tuples
\[\left(R, (1, \omega, \theta), I, (\alpha_1, \alpha_2, \alpha_3),\delta\right)\] where:
\begin{itemize}
    \item $R$ is a nondegenerate cubic ring over $T$,
    \item $(1, \omega, \theta)$ is an ordered basis of $R$,
    \item $I$ is an
$R$-ideal that is free of rank $3$ as a $T$-module,
\item $(\alpha_1,\alpha_2, \alpha_3)$ is an ordered $T$ basis of $I$,
\item $\delta$ is an
invertible element of $R \otimes \Frac(T)$ such that $I^2 \subset (\delta)$ and $N_S(I)^2 = N(\delta)$.
\end{itemize}
Two such tuples
$(R,(1, \omega, \theta), I,(\alpha_1, \alpha_2, \alpha_3), \delta)$ and $(R',(1, \omega', \theta'),I',(\alpha_1', \alpha_2', \alpha_3'),\delta')$ are considered equivalent if there exists
a ring isomorphism $\phi \colon R \to R'$ mapping $(1, \omega, \theta)$ to $(1, \omega', \theta')$ and $(\alpha_1, \alpha_2, \alpha_3)$ to $(\alpha_1', \alpha_2', \alpha_3')$, together with an element
$\kappa \in R' \otimes \Frac(T)$ such that $I' = \kappa\,\phi(I)$ and
$\delta' = \kappa^2 \phi(\delta)$. Then $G_T$ acts naturally on $\mathcal{S}_T$. 

\begin{theorem}[Bhargava]
With the above notation, there is a natural $G_T$-equivariant bijection between the nondegenerate elements in $V_T$ and $\mathcal{S}_T$. The set of nondegenerate orbits for the action of
$G_T$ on $V_T$ is in canonical bijection with equivalence classes of triples
$(R,I,\delta)$, where $R$ is a nondegenerate cubic ring over $T$, $I$ is an
$R$-ideal that is free of rank three as a $T$-module, and $\delta$ is an
invertible element of $R \otimes \Frac(T)$ satisfying the conditions that
$I^2 \subset (\delta)$ and $N_S(I)^2 = N(\delta)$. Two such triples
$(R,I,\delta)$ and $(R',I',\delta')$ are considered equivalent if there exists
a ring isomorphism $\phi \colon R \to R'$ together with an element
$\kappa \in R' \otimes \Frac(T)$ such that $I' = \kappa\,\phi(I)$ and
$\delta' = \kappa^2 \phi(\delta)$.
\end{theorem}

\begin{proof}
The above result is proven in \cite{bhargavahighercomplaws2} (see also \cite[Theorem~11]{BhargavaVarma}).
\end{proof}
\par Given a non-degenerate pair $(A,B)\in V_{\mathbb{R}}$, the associated conics intersect in $4$ points in $\mathbb{P}^1(\mathbb{C})$. The action of $G_{\mathbb{R}}$ on $V_{\mathbb{R}}$ consists of $3$ nondegenerate orbits $V_{\mathbb{R}}^{(i)}$, where $i=0,1,2$. Here, $V_{\mathbb{R}}^{(i)}$ is the the $G_{\mathbb{R}}$ orbit consisting of nondegerate pairs $(A,B)\in V_{\mathbb{R}}$ for which the points of intersection consist of $i$ pairs of complex zeroes and $4-2i$ real zeroes. A pair $(A,B)\in V_{\Z}$ is said to be absolutely irreducible if the following conditions are satisfied:
\begin{itemize}
    \item $A$ and $B$ do not possess a common $\mathbb{Q}$-rational zero as conics in $\mathbb{P}^{2}$; and
    \item the binary cubic form $f(x, y)=\operatorname{Det}(A x-B y)$ is irreducible over $\mathbb{Q}$.
\end{itemize}
We set $V_\Z^{(i)}:=V_{\mathbb{R}}^{(i)}\cap V_{\Z}$. Consider an element $(A,B)\in V_\Z^{(i)}$ for which the associated binary cubic form is irreducible and let $(R, I, \delta)$ be the associated triple. The if $i=0$ or $2$, then $R$ is an order in a totally real cubic field. On the other hand, if $i=1$ then $R$ is an order in a complex cubic field. If $i=0$ (resp. $i=2$) then $\delta$ is (resp. is not) a totally positive element of $R\otimes \Q$. We refer to \cite[Lemma 21]{BhargavaVarma} and its proof for further details.
\subsection{Class groups and narrow class groups}
\par Consider pairs $(R, I, \delta)$ where $R$ is an order in an $S_3$ cubic number field. We note that such cubic rings make up $100\%$ of all orders in cubic fields. A triple $(R,I,\delta)$ is called \emph{projective} if $I$ is a projective
$R$-module, equivalently an invertible fractional ideal; in this case
$I^{2}=(\delta)$. For a fixed $S_{3}$-cubic order $R$, equivalence classes of
projective triples admit a natural composition law
\[
(R,I,\delta)\circ (R,I',\delta')=(R,II',\delta\delta'),
\]
which makes them into a group, denoted $H(R)$.

Let $U$ be the unit group of $R$, $U^{\mathrm{pos\,norm}}$ the subgroup of units
of positive norm, and $U^{2}$ the subgroup of square units. There is an exact
sequence
\[
1 \longrightarrow \frac{U^{\mathrm{pos\,norm}}}{U^{2}}
\longrightarrow H(R)
\longrightarrow \mathrm{Cl}_{2}(R)
\longrightarrow 1.
\]
If $R$ is an order in a totally real (resp. complex) $S_{3}$-cubic field then $|U^{\mathrm{pos\,norm}}/U^{2}|=4$ (resp. $=2$). Thus we find that if $R$ is an order in a totally real (resp. complex) $S_{3}$-cubic field then
$|H(R)|=4\,|\mathrm{Cl}_{2}(R)|$ (resp. $|H(R)|=2\,|\mathrm{Cl}_{2}(R)|$).

When $R$ is totally real, let $H^{+}(R)\subset H(R)$ be the subgroup of triples
$(R,I,\delta)$ with $\delta$ totally positive. Then, one has the relation
\[
|H^{+}(R)|=|\mathrm{Cl}_{2}^{+}(R)|
\]
cf. \cite[Lemma 15]{BhargavaVarma}.

\section{Volume computations}
\par Let $W\subseteq \mathcal{S}_2$ be a measurable subset whose boundary has measure zero. Given a subset $S$ of $V_{\Z}^{(i)}$ and a positive real number $X$, we denote by $N(S,X)$ the number of $\op{G}_{\Z}$-orbits of irreducible elements $x\in S$ such that $|\op{Disc}(x)|<X$. Given an irreducible $x\in S$ let $\op{sh}(x)$ be the shape of the associated cubic ring. Denote by $N_W(S,X)$ the number of $\op{G}_{\Z}$-orbits of irreducible elements $x\in S$ such that $|\op{Disc}(x)|<X$ and such that $\op{sh}(x)\in W$. 
\par Denote by $\op{GL}_n^{\pm1}$ the subgroup of $\op{GL}_n$ with determinant $+1$ or $-1$ and set $\mathbb{G}_m:=\op{GL}_1$. The space $V_{\mathbb{R}}$ with its action by $G_{\mathbb{R}}$ is also considered by Bhargava and Harron \cite{bhargavaharron}. The key difference between our setting and theirs is that we use irreducible orbits in $V_{\Z}/G_{\Z}$ to parameterize $2$-torsion ideal classes in cubic rings, while in \emph{loc. cit.} they parameterize quartic rings. We are interested the distribution of the associated cubic field which is captured by the action of the $\op{GL}_2$ component of $G_{\mathbb{R}}$. On the other hand in \emph{loc. cit.} the shape of the quartic field is studied, which is encoded by the action of the $\op{GL}_3$ component of $G_{\mathbb{R}}$. Following Bhargava and Harron, the action of $G_{\mathbb{R}}$ on $V_{\mathbb{R}}$ factors through that of \[G_{\mathbb{R}}^{\prime}:=\mathbb{G}_{m}(\mathbb{R}) \times \mathrm{GL}_{3}^{ \pm 1}(\mathbb{R}) \times \mathrm{GL}_{2}^{ \pm 1}(\mathbb{R})\] via the homomorphism $G_{\mathbb{R}}\rightarrow G_{\mathbb{R}}'$ which sends $(g_3,g_2)\in \op{GL}_3(\mathbb{R})\times \op{GL}_2(\mathbb{R})$ to 
\[\left(\left|\operatorname{det} g_{3}\right|^{2 /3}\left|\operatorname{det} g_{2}\right|^{1 /2}, g_{3}^{\prime}, g_{2}^{\prime}\right)\]
\noindent where $g_{i}^{\prime}$ is given by $g_{i}=\left|\operatorname{det} g_{i}\right|^{1 / i} g_{i}^{\prime}$. The kernel of the action of $G_{\mathbb{R}}^{\prime}$ on $V_{\mathbb{R}}$ consists of $4$ elements. Fix a vector $v^{(i)}\in V_{\mathbb{R}}^{(i)}$ for which \[\op{sh}(v^{(i)})=\mtx{1}{0}{0}{1}\in \mathcal{S}_2\text{ and }|\op{Disc}(v^{(i)})|=1.\] Note that for $\gamma=(\lambda, g_3, g_2)\in G_{\mathbb{R}}'$, we have that 
\[\op{sh}(\gamma\cdot v^{(i)})=g_2\cdot \op{sh}(v^{(i)})=g_2.\]
Let $\mathcal{F} \subset G_{\mathbb{R}}^{\prime}$ be a fundamental domain for the left action of $G_{\mathbb{Z}}^{\prime}$ on $G_{\mathbb{R}}^{\prime}$ and set 
\[\begin{split}
    &\mathcal{R}_{X}:=\left\{x \in \mathcal{F} v^{(i)}\mid |\operatorname{Disc}(x)|<X\right\},\\
    &\mathcal{R}_{X, W}:=\left\{x \in \mathcal{F} v^{(i)}\mid |\operatorname{Disc}(x)|<X\text{ and }\op{sh}(x)\in W\right\}.
\end{split}\]
Letting $4 n_{i}$ be the cardinality of the stabilizer in $G_{\mathbb{R}}^{\prime}$ of $v^{(i)} \in V_{\mathbb{R}}^{(i)}$, one has that\[
n_0 = 24, \qquad n_1 = 4, \qquad n_2 = 8.
\]
\begin{theorem}
    The the following assertions hold:
    \begin{enumerate}
        \item $N\left(V_{\mathbb{Z}}^{(i)} ; X\right)=\frac{1}{n_{i}} \operatorname{Vol}\left(\mathcal{R}_{X}\right)+o(X)=\frac{1}{n_{i}} \operatorname{Vol}\left(\mathcal{R}_{1}\right) \cdot X+o(X)$,
        \item $N\left(V_{\mathbb{Z}}^{(i)} ; X, W\right)=\left(1 / n_{i}\right) \operatorname{Vol}\left(\mathcal{R}_{X, W}\right)+o(X)=\left(1 / n_{i}\right) \operatorname{Vol}\left(\mathcal{R}_{1, W}\right) \cdot X+o(X)$.
    \end{enumerate}
\end{theorem}
\begin{proof}
    The first assertion is \cite[Theorem 4]{bhargavaharron}. The proof of the second assertion is identical to that of Theorem 5 of \emph{loc. cit.} and is therefore omitted.
\end{proof}


\begin{theorem}\label{volume quotient}
     For $n \in\{3,4,5\}$, we have
$$
\frac{\operatorname{Vol}\left(\mathcal{R}_{1, W}\right)}{\operatorname{Vol}\left(\mathcal{R}_{1}\right)}=\frac{\mu(W)}{\mu\left(\mathcal{S}_{2}\right)} .
$$
\end{theorem}
\begin{proof}
     The proof of this result is similar to that of \cite[Theorem 7]{bhargavaharron}, the only difference lies in the fact that the shape is determined by the $\op{GL}_2$-action from the third coordinate of $G_{\mathbb{R}}'$ (and not the $\op{GL}_3$-action from the second coordinate). It follows from Proposition 12 of \emph{loc. cit.} that
\begin{equation*}
\frac{\operatorname{Vol}\left(\mathcal{R}_{1, W}\right)}{\operatorname{Vol}\left(\mathcal{R}_{1}\right)}=\frac{\int_{g \in G_{\mathbb{R}}^{\prime}} \chi_{\mathcal{R}_{1, W}}\left(g \cdot v^{(i)}\right)\left|\operatorname{Disc}\left(g \cdot v^{(i)}\right)\right| d g}{\int_{g \in G_{\mathbb{R}}^{\prime}} \chi_{\mathcal{R}_{1}}\left(g \cdot v^{(i)}\right)\left|\operatorname{Disc}\left(g \cdot v^{(i)}\right)\right| d g} .
\end{equation*}

If $g=\left(\lambda, g_3, g_2\right) \in G_{\mathbb{R}}^{\prime}$, then the first two coordinates $\lambda$ and $g_3$ do not affect the shape. Arguing as in \emph{loc. cit.}, we find that
$$
\frac{\operatorname{Vol}\left(\mathcal{R}_{1, W}\right)}{\operatorname{Vol}\left(\mathcal{R}_{1}\right)}=\frac{\int_{g \in \mathrm{GL}_{2}^{ \pm 1}(\mathbb{Z}) \backslash \mathrm{GL}_{2}^{ \pm 1}(\mathbb{R}) / \mathrm{GO}_{2}^{ \pm 1}(\mathbb{R})} \chi_{\mathcal{R}_{1, W}}\left(g \cdot v^{(i)}\right) d g \int_{k \in \mathrm{GO}_{2}^{ \pm 1}(\mathbb{R})} d k}{\int_{g \in \mathrm{GL}_{2}^{ \pm 1}(\mathbb{Z}) \backslash \mathrm{GL}_{2}^{ \pm 1}(\mathbb{R}) / \mathrm{GO}_{2}^{ \pm 1}(\mathbb{R})} d g \int_{k \in \mathrm{GO}_{2}^{ \pm 1}(\mathbb{R})} d k} =\frac{\mu(W)}{\mu\left(\mathcal{S}_{2}\right)}.
$$
\end{proof}

\section{Congruence conditions}

In this section we study local conditions at all primes $p$. We shall recall some results from \cite{BhargavaVarma} regarding the distribution of $2$-torsion ideals of cubic rings ordered according to discriminant. Let $S \subset V_{\mathbb{Z}}$ be a $G_{\mathbb{Z}}$-invariant subset defined by
finitely many congruence conditions, encoding prescribed local arithmetic
constraints such as maximality, splitting, or ramification. For each prime $p$,
its closure $S_p \subset V_{\mathbb{Z}_p}$ is a compact open subset stable under
the action of $G_{\mathbb{Z}_p}$. To $S_p$ one associates a local $p$-adic mass
\[
M_p(S)
=
\sum_{x \in G_{\mathbb{Z}_p}\backslash S_p}
\frac{1}{\operatorname{Disc}_p(x)\,|\operatorname{Stab}_{G_{\mathbb{Z}_p}}(x)|},
\]
which depends only on the local specification of $S$ at $p$; the resulting Euler
product converges since $S$ is defined by finitely many congruence conditions.

Let $N^{(i)}(S;X)$ denote the number of $G_{\mathbb{Z}}$-equivalence classes of
absolutely irreducible elements in $S \cap V_{\mathbb{Z}}^{(i)}$ with
$|\operatorname{Disc}(A,B)|<X$. Then \cite[Theorem 19]{BhargavaVarma} asserts that:
\[
\lim_{X\to\infty} \frac{N^{(i)}(S;X)}{X}
=
\frac{1}{2n_i}
\prod_p \left(\frac{p-1}{p}\, M_p(S)\right).
\]
Next consider sets defined by infinitely many congruence conditions. For each prime \(p\), let \(\Sigma_p\) be a set of isomorphism
classes of nondegenerate cubic rings over \(\mathbb{Z}_p\). 

\begin{definition}\label{acceptable Sigma}A collection
\(\Sigma=(\Sigma_p)_p\) is said to be \emph{acceptable} if, for all sufficiently
large primes \(p\), the set \(\Sigma_p\) contains every maximal cubic ring over
\(\mathbb{Z}_p\) that is not totally ramified.
\end{definition}
For a cubic order \(R\) over
\(\mathbb{Z}\), we write \(R\in\Sigma\), or say that \(R\) is a \(\Sigma\)-order,
if \(R\otimes\mathbb{Z}_p\in\Sigma_p\) for all primes \(p\). Fix such an acceptable collection \(\Sigma\), and let \(S=S(\Sigma)\) denote the
subset of absolutely irreducible elements \((A,B)\in V_{\mathbb{Z}}\) for which,
under Bhargava’s correspondence, the associated triple \((R,I,\delta)\) satisfies
\(R\in\Sigma\). Then one has that
\[
\begin{split}\lim_{X\to\infty}
\frac{N^{(i)}(S(\Sigma);X)}{X}
=&
\frac{1}{2n_i}
\cdot
\prod_p
\left(
\frac{p-1}{p}\cdot M_p(S(\Sigma))
\right),\\
=&\frac{2}{n_i}
\cdot
\prod_p
\left(
\frac{p-1}{p}
\sum_{R\in\Sigma_p}
\frac{1}{\op{Disc}_p(R)}
\cdot
\frac{1}{\lvert\op{Aut}(R)\rvert}
\right).
\end{split}
\]
cf. \cite[(12)]{BhargavaVarma}. Denote the resulting Euler product by
\[
c_\Sigma
:=
\prod_p
\left(
\frac{p-1}{p}
\sum_{R\in\Sigma_p}
\frac{1}{\op{Disc}_p(R)}
\cdot
\frac{1}{\lvert\op{Aut}(R)\rvert}
\right).
\]

\begin{lemma}\label{lemma 4.2}
The following assertions hold:
\begin{enumerate}
\item
The number of totally real \(\Sigma\)-orders \(\mathcal{O}\) with
\(\lvert\op{Disc}(\mathcal{O})\rvert<X\) is
\[
\frac{1}{12}\,c_\Sigma\,X + o(X).
\]
\item
The number of complex \(\Sigma\)-orders \(\mathcal{O}\) with
\(\lvert\op{Disc}(\mathcal{O})\rvert<X\) is
\[
\frac{1}{4}\,c_\Sigma\,X + o(X).
\]
\end{enumerate}
\end{lemma}
\begin{proof}
    The above result is \cite[Lemma 23]{BhargavaVarma}.
\end{proof}
\par Now let's fix a measurable set $W$ in $\mathcal{S}_2$ whose boundary has measure $0$ and let $N^{(i)}_W(S;X) $ be the number of $G_{\Z}$-equivalence classes of absolutely irreducible elements $(A,B)$ in $S\cap V_{\Z}^{(i)}$ for which $|\op{Disc}(A,B)|<X$ and $\op{sh}(A,B)\in W$.

\begin{lemma}\label{finitely many conditions}
   When $S$ is defined by finitely many local conditions, \[
\lim_{X\to\infty} \frac{N^{(i)}_W(S;X)}{X}
=
\frac{1}{2n_i}\frac{\mu(W)}{\mu(\mathcal{S}_2)}
\prod_p \left(\frac{p-1}{p}\, M_p(S)\right).
\]
\end{lemma}
\begin{proof}
    It suffices to show that 
    \[\lim_{X\rightarrow \infty} \frac{N^{(i)}_W(S;X)}{N^{(i)}(S;X)}=\frac{\mu(W)}{\mu(\mathcal{S}_2)}.\]
    Let $\mu_p(S)$ be the local density of $S$ at $p$. It follows from \cite[Lemma 9]{bhargavaharron} that
    \[\begin{split}
        \lim_{X\rightarrow \infty} \frac{N^{(i)}_W(S;X)}{N^{(i)}(S;X)}=\frac{\frac{1}{n_i}\prod_p \mu_p(S)\op{Vol}(\mathcal{R}_{X,W})}{\frac{1}{n_i}\prod_p \mu_p(S)\op{Vol}(\mathcal{R}_{X})}.
    \end{split}\]
    The result follows from Theorem \ref{volume quotient}, which implies that
    \[\frac{\op{Vol}(\mathcal{R}_{X,W})}{\op{Vol}(\mathcal{R}_{X})}=\frac{\op{Vol}(\mathcal{R}_{1,W})}{\op{Vol}(\mathcal{R}_{1})}=\frac{\mu(W)}{\mu(\mathcal{S}_2)}.\]
\end{proof}

\begin{lemma}\label{main lemma 1}
    Let $\Sigma$ be acceptable, then
    \[\frac{N^{(i)}_W(S(\Sigma);X)}{X}=\frac{c_\Sigma}{2n_i}\frac{\mu(W)}{\mu(\mathcal{S}_2)}.\]
\end{lemma}
\begin{proof}
    The proof follows from Lemma \ref{finitely many conditions} in which finitely many local conditions are considered and a standard sieve argument (see \cite[section 5]{bhargavaharron} or the proof of \cite[(12)]{BhargavaVarma}).  
\end{proof}

\begin{lemma}\label{main lemma 2}
The following assertions hold:
\begin{enumerate}
\item
The number of totally real cubic \(\Sigma\)-orders \(\mathcal{O}\) with shape in $W$ and
\(\lvert\op{Disc}(\mathcal{O})\rvert<X\) is is asymptotic to
\[
\frac{c_\Sigma}{12}\frac{\mu(W)}{\mu(\mathcal{S}_2)}X.
\]
\item
The number of complex cubic \(\Sigma\)-orders \(\mathcal{O}\) with shape in $W$ and
\(\lvert\op{Disc}(\mathcal{O})\rvert<X\) is asymptotic to
\[
\frac{c_\Sigma}{4}\frac{\mu(W)}{\mu(\mathcal{S}_2)} X.
\]
\end{enumerate}
\end{lemma}

\begin{proof}
    We prove part (1), the proof of (2) is similar. It follows the results in \cite[section 5]{bhargavaharron} that 
    \[\lim_{X\rightarrow \infty} \frac{\# \left\{\cO\mid \cO\text{ is a totally real cubic }\Sigma\text{-order, }\op{sh}(\cO)\in W\text{, } |\op{Disc}(\cO)|<X\right\}}{\# \left\{\cO\mid \cO\text{ is a totally real cubic }\Sigma\text{-order, } |\op{Disc}(\cO)|<X\right\}}=\frac{\op{Vol}(\mathcal{R}_{1,W})}{\op{Vol}(\mathcal{R}_{1})}=\frac{\mu(W)}{\mu(\mathcal{S}_2)}.\]
    Part (1) of Lemma \ref{lemma 4.2} asserts that 
    \[\# \left\{\cO\mid \cO\text{ is a totally real cubic }\Sigma\text{-order, } |\op{Disc}(\cO)|<X\right\}=\frac{1}{12}\,c_\Sigma\,X + o(X)\] from which the result follows.
\end{proof}

\begin{theorem}\label{main thm of article}
    For positive discriminant,
\begin{equation}
\label{eq:avg-positive}
\lim_{X\rightarrow \infty} \frac{
\sum_{\substack{R\in\Sigma \\ 0<\op{Disc}(R)<X\\
\op{sh}(R)\in W}}
\bigl(\lvert\op{Cl}_2(R)\rvert - \tfrac14 \lvert\mathcal{I}_2(R)\rvert\bigr)
}{
\sum_{\substack{R\in\Sigma \\ 0<\op{Disc}(R)<X\\
\op{sh}(R)\in W}} 1
}
=1.
\end{equation}
Similarly, for negative discriminant,
\begin{equation}
\label{eq:avg-negative}
\lim_{X\rightarrow \infty} \frac{
\sum_{\substack{R\in\Sigma \\ -X<\op{Disc}(R)<0 \\ \op{sh}(R)\in W}}
\bigl(\lvert\op{Cl}_2(R)\rvert - \tfrac12 \lvert\mathcal{I}_2(R)\rvert\bigr)
}{
\sum_{\substack{R\in\Sigma \\ -X<\op{Disc}(R)<0 \\ \op{sh}(R)\in W}} 1
}
=1.
\end{equation}
\noindent Finally, restricting to totally real rings, one has that:
\begin{equation}
\label{eq:avg-totally-real}
\lim_{X\rightarrow \infty} \frac{
\sum_{\substack{R\in\Sigma \\ 0<\op{Disc}(R)<X\\\op{sh}(R)\in W}}
\lvert\op{Cl}_2^{+}(R)\rvert
}{
\sum_{\substack{R\in\Sigma \\ 0<\op{Disc}(R)<X \\ \op{sh}(R)\in W}} 1
}
=1.
\end{equation}
\end{theorem}
\begin{proof}
We argue exactly as in \cite[(17)]{BhargavaVarma}, replacing their global set
$S(\Sigma)$ by the present set defined by the shape condition
$\op{sh}(R)\in W$, and invoking Lemmas~\ref{main lemma 1} and
\ref{main lemma 2} in place of their counting theorems. Concretely, one has
\[
N^{(0)}(S;X)+N^{(2)}(S;X)
=
\sum_{\substack{R\in\Sigma\\0<\op{Disc}(R)<X\\ \op{sh}(R)\in W}}
\bigl(4|\op{Cl}_2(R)|-|\mathcal{I}_2(R)|\bigr),
\]
\[
N^{(1)}(S;X)
=
\sum_{\substack{R\in\Sigma\\-X<\op{Disc}(R)<0\\ \op{sh}(R)\in W}}
\bigl(2|\op{Cl}_2(R)|-|\mathcal{I}_2(R)|\bigr),
\]
and
\[
N^{(0)}(S;X)
=
\sum_{\substack{R\in\Sigma\\0<\op{Disc}(R)<X\\ \op{sh}(R)\in W}}
\bigl(|\op{Cl}_2^{+}(R)|-|\mathcal{I}_2(R)|\bigr).
\]
\noindent Dividing each identity by
\[
\sum_{\substack{R\in\Sigma\\ \pm\op{Disc}(R)<X\\ \op{sh}(R)\in W}} 1
\]
and letting $X\to\infty$, the main counting lemmas imply that the left-hand sides
tend to~$1$. Rearranging the resulting expressions yields
\eqref{eq:avg-positive} and \eqref{eq:avg-negative}. Finally, restricting to
totally real rings and using that
$|H^{+}(R)|=|\op{Cl}_2^{+}(R)|$ for such $R$, the same argument applied to
$N^{(0)}(S;X)$ gives \eqref{eq:avg-totally-real}. This completes the proof.
\end{proof}

\begin{proof}[Proof of Theorem \ref{thm a}]
    In order to prove Theorem \ref{thm a}, let $\Sigma$ be the acceptable collection of maximal cubic orders. Theorem 17 and Corollary 18 of \cite{BhargavaVarma} together imply that for any maximal cubic order, $|\mathcal{I}_2(R)|=1$. The result then immediately follows from Theorem \ref{main thm of article}.
\end{proof}

\begin{proof}[Proof of Theorem \ref{thm c}]
    The proof follows immediately from Theorem \ref{thm a} and the proof of \cite[Corollary 3]{BhargavaVarma}.
\end{proof}

\bibliographystyle{alpha}
\bibliography{references}

\end{document}
